\chapter{涌现与框架形成}
\label{ch:emergence}

\section{为什么涌现重要}

较早版本的手稿说话的方式，仿佛一个框架已经预先存在，
而人所需要做的只是抵抗它或把它重定向。
这并不完整。
持久的框架通常是通过重复暴露、奖赏历史、环境耦合以及既往数据的干扰而形成的。
如果本书不解释框架是如何形成的，它就只能描述失败，而无法说明起源。

本章主要属于\textbf{中等证据}章节。
为本章建立索引的关键现代来源是
\emph{Metastability demystified: the foundational past, the pragmatic present and the promising future}
（Nature Reviews Neuroscience, 2024）以及
\emph{Associative conditioning in gene regulatory network models increases integrative causal emergence}
（Communications Biology, 2025）。
这些来源并不直接测量学习常规或个人习惯。
它们的价值在于机制构建：它们阐明了持久的高层组织如何能够从重复的局部互动中产生，
而不需要一个中央控制者。
任何把这些模型直接迁移到日常人类自我管理中的做法，仍然属于解释性推断，而不是强实证证明。

\section{框架作为历史产物}

\begin{definition}[框架形成]
\textbf{框架形成}是这样一个过程：重复输入、行为排练与环境规律共同生成一个相对稳定的行为或认知机制。
\end{definition}

\begin{proposition}[历史敏感的惯性]
当前状态的惯性不仅依赖当前刺激，还依赖强化、预期与成功复用的历史积累。
\end{proposition}

\begin{discussion}
这一命题很重要，因为它解释了为什么一个经过良好训练的框架常常会让人感觉“自然”。
它并不是那种一劳永逸被赋予的自然。
它之所以显得自然，是因为它被反复排练到让系统能够以很低成本进入其中。
\end{discussion}

\section{从局部互动到整体模式}

涌现补上了微观动作与宏观模式之间缺失的那一层。
重复的低层事件可以生成一个更高层次的结构，而后者随后又会反过来影响局部选择。
用本书的语言来说：
\begin{enumerate}[nosep]
  \item 重复行动会构建一个框架，
  \item 这个框架会改变后续行动的成本，
  \item 新的成本景观会稳定未来行为。
\end{enumerate}

\begin{example}[学习身份作为一种涌现机制]
一个学生反复在同一地点、用同一个进入线索、以同样的笔记本摆放方式开始工作，
并不只是在重复任务。
这个学生是在构建一个稳定机制，在其中学习状态会变得更容易再次进入。
这个框架并不是预先存在的。
它是由重复加上环境一致性涌现出来的。
\end{example}

项目中被索引的较早期涌现文献，尤其是 John
Holland 的 \emph{Emergence} 与 \emph{Introduction to the Theory of Complex
Systems}，为上述例子提供了教材层面的逻辑。
局部规则不会永远停留在局部。
一旦互动重复得足够久，宏观组织就会反过来作用于下一轮微观决策。
其业务含义很简单：如果一个组织或一个人持续地产生同一种模式，
正确的问题不只是“是谁选择了这个？”，
还应当是“是什么局部规则与什么历史在不断再生它？”

\section{亚稳态，而非僵硬或混乱的二选一}

Nature Reviews Neuroscience 关于亚稳态的文章之所以重要，
是因为它阻止了一种常见错误。
框架不应被想象成要么完全僵硬，要么完全流动。
被该综述以中等强度支持的主张是：系统能够在维持有组织模式的同时，仍保有重构能力。
这正是本手稿所需要的条件。
一个有用的框架必须足够持久，以降低再次进入的成本，
但又不能僵硬到在任务、情境或激励发生变化时把行动者困住。

这也是为什么亚稳态比装饰性的复杂性语言更重要。
它为本章提供了一种克制的说法，去表明一个机制可以同时具备稳定、脆弱和可重构这三种性质。
在行为解释层面，这能够说明为什么有些常规是高产的，而另一些则会变成牢笼。

\section{联结条件作用与因果涌现}

2025 年那篇关于基因调控网络模型中联结条件作用的 Communications Biology 论文，
在这里被用于支撑一个中等强度的观点。
它展示了一条明确路径：重复的局部条件作用如何增加整合性的因果涌现。
本手稿并不需要把那个模型逐项复制到人类心理学中。
它只需要吸收其中的机制教训：重复的局部耦合可以产生一种更高层次的组织，
而这一组织随后本身就具有解释价值。

这对框架形成十分关键。
框架并不只是若干彼此分离习惯的集合。
一旦足够多的局部线索、奖赏预期与行动序列开始相互强化，
更高层的机制就会开始支配什么事情显得轻松、显得显著、以及显得更可能发生。
从管理角度看，这种模式会变成一个真实的运行层。

\section{干扰、数据与路径依赖}

用户关于“既往数据干扰”的要求，可以被精确到足以进入教材写作的程度。
系统会携带先前成功、先前失败、先前奖赏与先前情境的痕迹。
这些痕迹会制造路径依赖：同样的当前线索，
会因为系统此前已经学会期待什么，而产出不同的行动。

\begin{remark}
这正是涌现使手稿不至于过度个人主义化的地方。
框架会被既往数据、社会环境、习得的常规以及与特定奖赏结构的重复接触所塑形。
\end{remark}

这里必须保持证据分层。
路径依赖作为一般系统观念，受到本章中等证据来源与教材的良好支持。
把这一观念应用于个人的框架惯性，也同样是合理的。
但是，如果主张某一个具体的人类常规是由与基因调控模型完全相同的机制造成的，
那就会变成\textbf{推测性}内容。
因此，本章把类比保持在恰当的抽象层级：重复耦合会生成历史敏感的组织，而这种组织会改变后续选择。

\section{为什么形成过程对干预很重要}

一旦把框架形成理解为一个具有历史性的涌现过程，本书的干预逻辑就会得到改善。
只在失败发生的当下去打败一个坏机制，成本通常很高，
因为这个机制早已由许多更早的循环构建出来。
更便宜的策略，是重新设计那些不断生成该机制的重复性局部条件：
\begin{enumerate}[nosep]
  \item 稳定你想要的进入线索；
  \item 降低第一个动作的歧义；
  \item 让奖赏足够早地到来，以强化复用；
  \item 移除那些不断重建旧机制的重复性环境耦合。
\end{enumerate}

这就是涌现语言的业务价值。
它把本书从一种关于“抵抗”的理论，转变为一种关于“生产”的理论：
框架是被制造出来的，而凡是被制造出来的东西，都可以通过改变不断生产它的循环而被重新制造。

\section{推荐阅读}

\begin{readingbox}
\textbf{基础}
\begin{itemize}[nosep]
  \item John Holland, \emph{Emergence}.
  \item \emph{Introduction to the Theory of Complex Systems}.
\end{itemize}
\textbf{现代研究}
\begin{itemize}[nosep]
  \item \emph{Metastability demystified: the foundational past, the pragmatic
        present and the promising future} （Nature Reviews Neuroscience, 2024,
        medium）。
  \item \emph{Associative conditioning in gene regulatory network models
        increases integrative causal emergence} （Communications Biology, 2025,
        medium）。
\end{itemize}
\textbf{前沿}
\begin{itemize}[nosep]
  \item 只有在读者已经理解历史强化与路径依赖之后，才应引入前沿的涌现方案。
\end{itemize}
\end{readingbox}

\begin{summarybox}
本章解释了框架从何而来。
它的有据可依之处主要在于机制构建，而非直接测量：
重复的局部互动可以生成更高层的组织，亚稳态解释了为什么有用的框架既持久又可修订，
而因果涌现语言则说明为什么一个成熟机制可以成为真正的解释层。
这使手稿从一个静态的“抵抗模型”，转变为一个关于形成、巩固与路径依赖的发展模型。
\end{summarybox}